\documentclass[quick,pro]{debate}
\motion{我厌蠢，我错了 / 我厌蠢，我没错}
\stance{正方}{我错了}
\begin{document}
\debatetitle
\begin{thesisbox}常人的厌蠢，平常凭比较就判人蠢，尺子握在已经会了的人手里，自己看不见；就算看准了，劝了还犯，也会从“说了那么多遍”走到“他就是那么蠢”，回头不再讲理。这份厌平常的样子，该改。\end{thesisbox}

\section{定义与判准（标准表述）}
\begin{fields}
\field{厌蠢}{厌蠢，是觉得别人犯蠢时生出的嫌恶、反感，常带着想躲开的心；它可以只冲着那件事，也可以冲着这个人、看低他；它是一种态度，不等于骂人或欺负人。}
\field{蠢}{蠢，是在一般人看来本该想得到、做得对的事上，明显想错、做错；可能出于不会，也可能出于不肯，原因不在定义里预设。（一般人：了解他处境、讲道理的一般人。我：一个普通的、会说出这句话的我。）}
\field{判准}{判断这份厌对不对，就看一个清醒的我问完三句：看准了吗？对他公道吗？划得来吗？问完觉得这份厌平常的样子就该改，就是错了；站得住，只在个别场合调调分寸，就是没错。\sub{我方证}{平常凭比较判，看不准；就算看准了，也会从要求走到判词，越过尊重。}\sub{对方证}{平常站得住，毛病只在个别场合。}\sub{门槛}{平常就这样、只是改得掉，照样算错了。}\sub{越线}{这回先不说、一时走开，不越线；回头他开口先认定不值得听，才越线。}\sub{我方让出}{教过、被指出、又做错的看准；走神手滑的看准；对那件事的恼火；一时走开。}}
\field{对方提偏向定义或判准时的 30 秒口径}{“是厌不是鄙视”：两头都写了、都用“可以”，平常落在哪头是今天要辩的。“对得上就没错”：只回答了一问，没回答尊重和划得来。“我有权远离”：有权不等于没错。“情绪无对错”：无所谓对错不等于没错。}
\end{fields}

\section{我方论点}
\begin{tabularx}{\linewidth}{@{}L{4.2em}L{8.6em}Y L{6.6em}@{}}\toprule
\thead{标签}&\thead{一句话}&\thead{核心数据 / 学理 / 案例}&\thead{出处}\\\midrule
\textbf{拿我当一般人}（准）&平常凭比较判，不凭记录；尺子在会了的人手里，只往“他本该会”偏，自己看不见&后见之明偏差。《Hindsight ≠ foresight》：3 项实验，知情者高估不知情者本来能知道多少，largely unaware。Hom：招募 193 人、排除 10 人，替陌生同龄人估字谜，偏差 essentially the same（须认被要求设想同伴经历相同）。案例：教人坐高铁的评论区；新京报写同学同事&Fischhoff 1975；Hom 2023；上游新闻 2023；新京报 2023\\
\textbf{说了也没用}（公·尊重）&就算看准了，劝了还犯，会从“说了那么多遍”走到“他就是那么蠢”，回头不再讲理&《Beat them or ban them》：3 项研究，愤怒长期和解、轻蔑长短期排斥；轻蔑 may 从先前愤怒发展；管不了、品性归因预示轻蔑（须认 often occur together）。轻蔑贬品性多于行为。案例：劝了多次，一度降至冰点，现很微妙；随后单列“说了那么多遍，怎么还是那么蠢”（未标说话人）&Fischer \& Roseman 2007；Fischer \& Giner-Sorolla 2016；每日商报 2024\\
\bottomrule\end{tabularx}
\textbf{其余材料}：Weiner 等 1988《An attributional analysis of reactions to stigmas》2 项实验；Heslin 等 2006《Keen to help?》3 项研究；Wood \& Porter 2019《The Elusive Backfire Effect》5 个实验、10,100 余人、52 个议题；Edmondson 1999《Psychological Safety and Learning Behavior in Work Teams》51 个团队，相关；Guilbault 等 2004 后见之明元分析，95 项研究，Md=.39，减偏操纵效应量不更低；Birch 等 2017，3 项实验、359 名本科生；Roese \& Vohs 2012《Hindsight bias》；康德 1797《道德形而上学》§39（自译，整句引）；Strawson 1962《自由与怨恨》。

\section{对方论点 → 一句话反驳}
\begin{tabularx}{\linewidth}{@{}L{9em}Y L{12em}@{}}\toprule
\thead{对方会说}&\thead{我方一句话回}&\thead{对方追问时}\\\midrule
教过了还犯，熟人记录在手&凭记录的我们让；第一次那一下只有比较；看准之后往哪走是第二点&“一问就知道”→ 那就是先问，请证常人平常问\\
情绪按能不能控制分流（64 项研究元分析）&只管判断之后，不管判断准不准；没线索时人默认读成不肯&“那是污名研究”→ 迁移是推论；元分析对象也不是蠢\\
恼他不弃他，恼里带着讲理&“不纯”不等于“没有”：常常一起出现，轻蔑可能从生气里长出来&“may 只是可能”→ 可能的条件正是劝了还犯\\
这回放下下回再来（孔子、莎莉文）&线同意；样子是老师和圣人，评的是常人&“常人也有”→ 也有就是个别\\
\bottomrule\end{tabularx}

\section{必问与必答}
\begin{points}{必问（不答就点名、重复、不换题）}
\item 常人说“我厌蠢”的那一刻，说出口的是“你卡在哪”，还是“这都不会”？（全场问题）
\item 没有记录时，说他本该会，凭的是不是比较？\quad 看准了他这件事，就能看低他这个人吗？
\end{points}
\begin{points}{必答（15 秒正面回答）}
\item 屡教不改怎么办？→ 对那件事生气、批评都可以，一时走开也可以；要改的是看低他，把他划到不用再讲理那边。
\item 回头不讲理是推论吗？→ 是，报道没写；依据是轻蔑不看好对方会变好，认定人改不了人就少教。
\item 只在心里厌算不算？→ 算，心里那份就是厌蠢本身，错不错看判准。要我们别生气？→ 不是，恼照样恼，先问再判。
\end{points}

\section{自由辩战场概要}
\begin{tabularx}{\linewidth}{@{}L{7em}Y Y@{}}\toprule
\thead{战场}&\thead{开场问题}&\thead{收束语}\\\midrule
凭什么判（准）&没有记录时，说他本该会，凭的是什么？&没有记录的那一下，凭的就是比较\\
劝了还犯之后（尊重）&看准了他这件事蠢，就能看低他这个人吗？&怪那件事我们不拦；要改的是走到判词那一截\\
平常还是个别（门槛）&平常就这样、只是改得掉，算不算错？&要学才有，就是平常没有\\
\bottomrule\end{tabularx}
前 90 秒战场一（抛完必问），中 90 秒战场二，后 60 秒战场三并收束，留 20 到 30 秒；一辩开场，四辩收束。

\section{如果……就……}
\begin{contingency}
\ifthen{对方拿出没预料到的数据}{先问机构、年份、报告名，再问测的是判断之后还是判断本身；报厌蠢比例的请出示来源。}
\ifthen{论点一被打成“只剩陌生人”}{收缩到论点二，从“就算看准了”起步，用门槛规则解释。}
\ifthen{被问“是不是推论”}{直接认，再给依据。对方回避必问：点名、重复、不换题；把“看低”写进定义：指出换题。}
\end{contingency}

各环节：正一立论 180 秒 → 反二质询 120 秒双边 → 反一立论 180 秒 → 正二质询 120 秒双边 → 反二驳论 120 秒 → 正二驳论 120 秒 → 四辩对辩各 90 秒 → 正三盘问 90 秒 → 反三盘问 90 秒 → 正三小结 120 秒 → 反三小结 120 秒 → 自由辩各 240 秒 → 反四结辩 210 秒 → 正四结辩 210 秒；奇袭质询（正三）或申论（正一）每方一次。

\begin{recordcard}
\record{反一立论结束}{对方立论的主干}{\opt{对方立论建立在「教过了还犯」上}\opt{对方立论建立在「恼他不弃他」上}\opt{对方立论建立在「情绪按能不能控制分流」上}\opt{对方立论建立在「情绪不经选择，无所谓对错」上}\opt{对方立论建立在「守标准、止损」的好处上}\opt{对方立论的主干不在以上几条时}}{正二驳论·驳对方立论（选 2）}
\record{反二驳论结束}{反二打了哪里}{\opt{反二打的是我方实验估的都是陌生人}\opt{反二打的是回头不讲理是推论}\opt{反二打的是先问本来就在反方的厌里}\opt{反二打的是我方只证明了有，没证明平常}\opt{反二没有打到我方论点}}{正二驳论·回应反二}
\record{盘问结束}{链一的结果}{\opt{对方说没有记录时常人也是先问}\opt{对方承认没有记录时常人是先判}\opt{反二、反四答得不一样}\opt{对方回避了链一的问题}\opt{链一没有问完}}{正三小结·链一}
\record{盘问结束}{链二的结果}{\opt{对方承认回头先认定不值得听算越线}\opt{对方说「那么蠢」说的是事}\opt{对方回避了链二的问题}\opt{链二没有问完}}{正三小结·链二}
\record{发动奇袭时}{在哪个窗口发动}{\opt{在盘问小结之后发动}\opt{在自由辩之后发动}\opt{发动时机不在以上几个窗口时}}{奇袭申论·为什么现在发动}
\record{自由辩结束}{全场主战场落在}{\opt{全场主战场落在「平常凭比较还是凭记录」}\opt{全场主战场落在「看准之后还讲不讲理」}\opt{全场主战场落在「先问是不是常人平常就有」}\opt{主战场不在以上几条时}}{正四结辩·归纳分歧}
\record{自由辩结束}{对方全场最有力的一点}{\opt{对方全场最有力的是「熟人之间记录在手」}\opt{对方全场最有力的是「回头不讲理是推论」}\opt{对方全场最有力的是「恼里带着讲理，不是轻蔑」}\opt{对方全场最有力的是「存在更好的，不等于这个就错了」}\opt{对方最有力的一点不在以上几条时}}{正四结辩·处理最强点}
\record{反四结辩结束}{反四收在哪里}{\opt{反四收在「你本可以」的恼是看得起他}\opt{反四收在「只在个别场合冒出来，那叫调分寸」}\opt{反四收在「存在更好的，不等于这个就错了」}\opt{反四收在「先问、讲理本来就在我们的厌里」}\opt{反四的收束不在以上几条时}}{正四结辩·回应反四（再接 60 字临场）}
\record{全场}{对方回避了哪些问题}{（写下来）}{小结与结辩里点名}
\end{recordcard}
\end{document}
